% Partie : logique du premier ordre
% Chapitre : syntaxe et sémantique
% Section : introduction

\documentclass[../../../include/open-logic-section]{subfiles}

\begin{document}

\olfileid[fr]{fol}{syn}{int}

\olsection{Introduction}

Pour développer la théorie et la métathéorie de la logique du
premier ordre, il nous faut d'abord définir la syntaxe et la
sémantique de ses expressions. Celles-ci sont les termes et
les !!{formula}s. Les termes sont construits à partir de
!!{variable}s, de !!{constant}s et de !!{function}s.
La combinaison de !!{predicate}s et de termes fournit les
!!{formula}s les plus simples, dites « atomiques ». À partir de ces
!!{formula}s atomiques, nous pouvons ensuite construire des
!!{formula}s plus complexes au moyen de connecteurs logiques
et de quantificateurs. Il existe de nombreuses façons de
définir les règles de formation ; nous n'en présentons qu'une.
D'autres systèmes choisiront d'autres symboles, d'autres
ensembles de connecteurs primitifs, ou emploieront autrement
les parenthèses, voire les supprimeront entièrement, comme
dans la notation dite polonaise. Toutes ces approches ont
cependant un point commun : les règles de formation
définissent \emph{inductivement} l'ensemble des termes et des
!!{formula}s. Lorsque ces règles sont convenablement
définies, chaque expression ne peut être construite que
d'une seule façon, à des différences inessentielles près.
Cette définition inductive donne des expressions
\emph{à lecture unique} et permet de leur attribuer un sens
par la même méthode : une définition inductive.

Attribuer un sens aux expressions relève de la sémantique.
Sa notion centrale est celle de satisfaction dans
!!a{structure}. Une !!{structure} attribue un sens aux
éléments constitutifs du langage. Son !!{domain} est un
ensemble non vide d'objets. Les quantificateurs portent sur
ce domaine ; les !!{constant}s sont interprétées par des
éléments du domaine, les !!{function}s par des fonctions de
l'arité correspondante, dont les arguments et les valeurs sont
dans ce domaine, et les !!{predicate}s par des relations sur le domaine.
Le !!{domain}, muni de ces interprétations du vocabulaire de
base, constitue !!a{structure}. Des !!{variable}s peuvent
figurer dans les !!{formula}s. Pour en donner la sémantique,
il faut aussi leur associer des !!{element}s du !!{domain} :
c'est ce que fait une assignation de variables.
La relation de satisfaction rassemble enfin ces éléments.
Une !!{formula} peut être satisfaite dans
!!a{structure}~$\Struct{M}$ relativement à une assignation de
variables~$s$ ; on écrit alors $\Sat{M}{!A}[s]$.
Cette relation est elle aussi définie par induction sur la
construction de~$!A$. Les tables de vérité des connecteurs
logiques permettent, par exemple, de définir la satisfaction
de $!A\land !B$ à partir de la satisfaction, ou de la
non-satisfaction, de $!A$ et de~$!B$.
Il s'avère alors que l'assignation de variables ne joue
aucun rôle lorsque la !!{formula}~$!A$ est !!a{sentence},
c'est-à-dire lorsqu'elle n'a aucune variable libre.
Nous pouvons donc parler simplement d'!!{sentence}s
satisfaits, ou non, dans des !!{structure}s.

À partir de la relation de satisfaction $\Sat{M}{!A}$ pour
les !!{sentence}s, nous pouvons définir les notions sémantiques
fondamentales de validité, de conséquence logique et de
satisfaisabilité. Un !!{sentence} est valide, ce que l'on note
$\Entails !A$, si toute structure le satisfait.
Il est une conséquence logique d'un ensemble d'%
!!{sentence}s, ce que l'on note $\Gamma\Entails !A$, si toute
!!{structure} qui satisfait tous les !!{sentence}s
de~$\Gamma$ satisfait aussi~$!A$.
Enfin, un ensemble d'!!{sentence}s est satisfaisable s'il existe
!!a{structure} qui satisfait simultanément tous ses
!!{sentence}s. Les !!{formula}s étant définies inductivement,
et la satisfaction étant à son tour définie par induction
sur leur construction, nous pouvons recourir à l'induction
pour établir des propriétés de notre sémantique et relier
les notions sémantiques ainsi définies.

\end{document}
